\documentclass[10pt,a4paper, english]{article}


\usepackage{graphicx}
\graphicspath{ {./images/} }

\usepackage{amssymb}
\usepackage{amsmath}


\usepackage{hyperref}

% hyperlink colors
\hypersetup{
         colorlinks,
         citecolor=black,
         filecolor=black,
         linkcolor=black,
         urlcolor=black,
         pdfborder=0 0 0
}

\usepackage{xcolor}

\usepackage[framemethod=TikZ]{mdframed}

\usepackage[left=1cm,right=1cm,top=2cm,bottom=1.5cm]{geometry}
\usepackage[skip=2pt]{caption} % example skip set to 2pt captions close to table figure etc

\usepackage{fancyhdr}

\usepackage{multicol}
\usepackage{float}

\usepackage{enumitem}
\setlist{topsep=0.4em, itemsep=0.1em}

\mdfsetup{% 
backgroundcolor=gray!20, 
hidealllines=true,
roundcorner=2.5pt}

\author{Daniel Schafhäutle}

\fancyhf{}

\fancyhead[C]{Explainable AI CDS120}

\fancyfoot[L]{\today}
\fancyfoot[C]{Daniel Schafhäutle}
\fancyfoot[R]{\thepage}

\pagestyle{fancy}


% silbentrennung stärke
\usepackage{microtype}
\usepackage{mwe}
\hyphenpenalty=-9000
\exhyphenpenalty=-9000





% ===============================
% Some Math Shortcuts
% ===============================
%\boldone for a \mathbb 1 (since numbers are not supported) Indicator
\newcommand*{\boldone}{\text{\usefont{U}{bbold}{m}{n}1}}

\providecommand{\iprod}[2]{\ensuremath{\left\langle #1,\,#2  \right\rangle}}
\providecommand{\mnorm}[1]{\ensuremath{\left\lvert#1\right\rvert}}
\providecommand{\norm}[1]{\ensuremath{\left\lVert#1\right\rVert }}


%some math shortcuts
\newcommand{\argmax}[1]{\mathop{\arg\max}\limits_{#1}}
\newcommand{\argmin}[1]{\mathop{\arg\min}\limits_{#1}}
\newcommand{\maxover}[1]{\mathop{\max}\limits_{#1}}
\newcommand{\minover}[1]{\mathop{\min}\limits_{#1}}
\newcommand{\infover}[1]{\mathop{\inf}\limits_{#1}}

% Define \E with an optional subscript and mandatory content
\newcommand{\E}[1]{\mathbb{E}\left[ #1 \right]}

% Expecation with respect to
\newcommand{\Ewrp}[2]{\mathbb{E}_{#1}\left[ #2 \right]}



% Left-right bracket
\newcommand{\lr}[1]{\left (#1\right)}

% Left-right square bracket
\newcommand{\lrs}[1]{\left [#1 \right]}

% Left-right curly bracket
\newcommand{\lrc}[1]{\left \{#1\right\}}

% Left-right absolute value
\newcommand{\lra}[1]{\left |#1\right|}

% Left-right upper value
\newcommand{\lru}[1]{\left \lceil#1\right\rceil}

% Scalar product
\newcommand{\vp}[2]{\left \langle #1 , #2 \right \rangle}

% The real numbers
\newcommand{\R}{\mathbb R}

% The natural numbers
\newcommand{\N}{\mathbb N}



% Indicator function with an optional argument
\NewDocumentCommand{\1}{o}{\mathds 1{\IfValueT{#1}{\lr{#1}}}}

% Probability function
\let\P\undefined
\NewDocumentCommand{\P}{o}{\mathbb P{\IfValueT{#1}{\lr{#1}}}}

% A hypothesis space
\newcommand{\HH}{\mathcal H}

% A sample space
\newcommand{\XX}{\mathcal{X}}

% A label space
\newcommand{\YY}{\mathcal{Y}}

% A nicer emptyset symbol
\let\emptyset\varnothing

% Sign operator
\DeclareMathOperator{\sign}{sign}
\newcommand{\sgn}[1]{\sign\lr{#1}}

% KL operator
\DeclareMathOperator{\KL}{KL}

% kl operator
\DeclareMathOperator{\kl}{kl}

% The entropy
\let\H\relax
\DeclareMathOperator{\H}{H}

% Majority vote
\DeclareMathOperator{\MV}{MV}

% Variance
\DeclareMathOperator{\V}{Var}
\NewDocumentCommand{\Var}{o}{\V\IfValueT{#1}{\lrs{#1}}}

% VC
\DeclareMathOperator{\VC}{VC}

% VC-dimension
\newcommand{\dVC}{d_{\VC}}

% FAT ...
\DeclareMathOperator{\FAT}{FAT}
\newcommand{\dfat}{d_{\FAT}}
\newcommand{\lfat}{\ell_{\FAT}}
\newcommand{\Lfat}{L_{\FAT}}
\newcommand{\hatLfat}{\hat L_{\FAT}}

% Distance
\DeclareMathOperator{\dist}{dist}




%%%%%%%%%%%%%%%%%%%%%%%
% Reference commands  %
%%%%%%%%%%%%%%%%%%%%%%%
\newcommand{\Figref}[1]{Figure~\ref{#1}}  % Figure reference
\newcommand{\Secref}[1]{Section~\ref{#1}}  % Section reference
\newcommand{\Eqref}[1]{Equation~(\ref{#1})}  % Equation reference
\newcommand{\Algoref}[1]{Algorithm~\ref{#1}}  % Algorithm reference
\newcommand{\Defref}[1]{Definition~\ref{#1}}    % Definition reference
\newcommand{\Thmref}[1]{Theorem~\ref{#1}}   % Theorem reference
\newcommand{\Lemref}[1]{Lemma~\ref{#1}}   % Lemma reference
\newcommand{\Cororef}[1]{Corollary~\ref{#1}}    % Corollary reference
\newcommand{\Propref}[1]{Proposition~\ref{#1}}  % Proposition reference


%%%%%%%%%%%%%%%%%%%%%%%%%%%%%%%%%%%%
% Definition, Theorems, proof etc. %
%%%%%%%%%%%%%%%%%%%%%%%%%%%%%%%%%%%%
\newtheorem{definition}{Definition}[section]  % [section]: recount after each section
\newtheorem{theorem}{Theorem}[section]
\newtheorem{lemma}[theorem]{Lemma}            % [theorem]: same counter as theorem
\newtheorem{corollary}[theorem]{Corollary}
\newtheorem{proposition}[theorem]{Proposition}
\newtheorem{hypothesis}[theorem]{Hypothesis}
\newtheorem{conjecture}[theorem]{Conjecture}
\newtheorem{principle}[theorem]{Principle}
\newtheorem{claim}[theorem]{Claim}
\newtheorem{example}[theorem]{Example}
% \begin{proof}\end{proof} already included in package amsthm

\newcommand{\KLDiv}[2]{D_{\mathrm{KL}} \left( {#1} \middle\| {#2} \right)}

\newcommand{\avgsum}[3]{\frac{1}{ {#3} } \sum_{ {#1}={#2} }^{ {#3} }}


% =================
% CODE
% =================

\usepackage{xcolor} % For a colorfull presentation
\usepackage{listings} % For presenting code
\usepackage{algorithm}
\usepackage{algpseudocode}
\usepackage{subcaption}
% Definition of a style for code, matter of taste
\lstdefinestyle{mystyle}{
language=Python,
basicstyle=\ttfamily\footnotesize,
backgroundcolor=\color[HTML]{F7F7F7},
rulecolor=\color[HTML]{EEEEEE},
identifierstyle=\color[HTML]{24292E},
emphstyle=\color[HTML]{005CC5},
keywordstyle=\color[HTML]{D73A49},
commentstyle=\color[HTML]{6A737D},
stringstyle=\color[HTML]{032F62},
emph={@property,self,range,True,False},
morekeywords={super,with,as,lambda},
literate=%
{+}{{{\color[HTML]{D73A49}+}}}1
{-}{{{\color[HTML]{D73A49}-}}}1
{*}{{{\color[HTML]{D73A49}*}}}1
{/}{{{\color[HTML]{D73A49}/}}}1
{=}{{{\color[HTML]{D73A49}=}}}1
{/=}{{{\color[HTML]{D73A49}=}}}1,
breakatwhitespace=false,
breaklines=true,
captionpos=b,
keepspaces=true,
numbers=none,
showspaces=false,
showstringspaces=false,
showtabs=false,
tabsize=4,
frame=single,
}
\lstset{style=mystyle}

\usepackage{minted}
\usemintedstyle{colorful}
\usepackage{xcolor} % to access the named colour LightGray
\definecolor{LightGray}{gray}{0.95}

\usepackage{hyperref}

% hyperlink colors
\hypersetup{
	colorlinks,
	citecolor=black,
	filecolor=black,
	linkcolor=black,
	urlcolor=black,
	pdfborder=0 0 0
}

\usepackage{xcolor}

\usepackage[framemethod=TikZ]{mdframed}

\usepackage[left=1cm,right=1cm,top=2cm,bottom=1.5cm]{geometry}
\usepackage[skip=2pt]{caption} % example skip set to 2pt captions close to table figure etc

\usepackage{fancyhdr}

\usepackage{multicol}
\usepackage{float}

\usepackage{enumitem}
\setlist{topsep=0.4em, itemsep=0.1em}

\mdfsetup{% 
	backgroundcolor=gray!20,
	hidealllines=true,
	roundcorner=2.5pt}

\author{Daniel Schafhäutle}

\fancyhf{}

\fancyhead[C]{Explainable AI CDS120}

\fancyfoot[L]{\today}
\fancyfoot[C]{Daniel Schafhäutle}
\fancyfoot[R]{\thepage}

\pagestyle{fancy}


% silbentrennung stärke
\usepackage{microtype}
\usepackage{mwe}
\hyphenpenalty=-9000
\exhyphenpenalty=-9000





% =================
% CODE
% =================

\usepackage{xcolor} % For a colorfull presentation
\usepackage{listings} % For presenting code
\usepackage{algorithm}
\usepackage{algpseudocode}
\usepackage{subcaption}
% Definition of a style for code, matter of taste
\lstdefinestyle{mystyle}{
language=Python,
basicstyle=\ttfamily\footnotesize,
backgroundcolor=\color[HTML]{F7F7F7},
rulecolor=\color[HTML]{EEEEEE},
identifierstyle=\color[HTML]{24292E},
emphstyle=\color[HTML]{005CC5},
keywordstyle=\color[HTML]{D73A49},
commentstyle=\color[HTML]{6A737D},
stringstyle=\color[HTML]{032F62},
emph={@property,self,range,True,False},
morekeywords={super,with,as,lambda},
literate=%
	{+}{{{\color[HTML]{D73A49}+}}}1
{-}{{{\color[HTML]{D73A49}-}}}1
{*}{{{\color[HTML]{D73A49}*}}}1
{/}{{{\color[HTML]{D73A49}/}}}1
{=}{{{\color[HTML]{D73A49}=}}}1
{/=}{{{\color[HTML]{D73A49}=}}}1,
breakatwhitespace=false,
breaklines=true,
captionpos=b,
keepspaces=true,
numbers=none,
showspaces=false,
showstringspaces=false,
showtabs=false,
tabsize=4,
frame=single,
}
\lstset{style=mystyle}

\usepackage{minted}
\usemintedstyle{colorful}
\usepackage{xcolor} % to access the named colour LightGray
\definecolor{LightGray}{gray}{0.95}

\begin{document}
\tableofcontents
\twocolumn

\section{Introduction to Explainable AI}
\subsection{Taxonomy}
By Interpretability Type
\begin{itemize}
	\item Intrinsic Interpretability
	\item Post-hoc Interpretability
\end{itemize}
By Result Type
\begin{itemize}
	\item Feature summary statistics
	\item Feature summary visualizations
	\item Model internals (weights, tree structure)
	\item Data Point base (counterfactuals, prototypes)
	\item Intrinsic surrogates
\end{itemize}
\subsection{Global vs Local interpretability}
\paragraph{Global Interpretability}
How well the trainied model makes predictions overall.
Which features are important for predictions overall.
Feature interactions.
\paragraph{Local Interpretability}
Why the model made a particular prediction for one instance
Locally a complex model may behave in a simple way.

\subsection{Making explanations}
When making explanations, we need to keep the human factor in mind.
They should be
\begin{itemize}
	\item Contrastive: Provide “why this, not that?” with a clear reference point.
	\item Selected: Limit to the top 1–3 drivers.
	\item Social: Match audience knowledge and context.
	\item Abnormality: Highlight unusual but influential factors.
	\item Fidelity: Stay faithful to model behavior (and reality when possible).
	\item Beliefs: Acknowledge/prioritize known priors or add constraints if needed.
	\item Generality: Indicate coverage/support when relevant.
\end{itemize}

\section{Interpretable Models}
Using inherently interpretable algorithms.
\subsection{Linear Regression}
Predicts target as weighted sum of inputs.
\begin{equation}
	y = \beta_0 + \beta_1 x_1 + \dots + \beta_p x_p + \epsilon
	\label{eq:linear_regression}
\end{equation}
Estimated usually by Ordinary Least Squares (OLS)
\begin{equation}
	\hat \beta = \argmin{\beta_0, \dots , \beta_p} \sum_{i=1}^n \lr{y^{(i)} - \lr{\beta_0 + \sum_{j=1}^p}\beta_jx_j^{(i)}}²
	\label{eq:OLS}
\end{equation}
Each weight represents the effect of one feature, holding others constant.
We can compute \textbf{confidence intervals} for each coefficient.
\paragraph{Regression Output}
\begin{table}[H]
	\caption{Regression Output}\label{tab:regression_output}
	\begin{center}
		\begin{tabular}{l|r|r|r}
			\hline
			Feature     & Weight$\hat \beta_j$ & SE  & $|t|$ \\
			\hline
			Temperature & 110.7                & 7.0 & 15.7  \\
			Humidity    & -17.4                & 3.2 & 5.5   \\
			\hline
		\end{tabular}
	\end{center}
\end{table}
\begin{itemize}
	\item Weight: Estimated influence of $x_i$ over the outcome $y$
	\item SE (Standard Error): Uncertainty (variance) of the estimated coefficient
	\item $|t|$ t-statistic: Strength of evidence that $\beta_j \neq 0$
\end{itemize}
\begin{equation}
	t_{\beta_j} = \frac{\hat \beta_j}{SE(\hat \beta_j)}
	\label{eq:t-statistic}
\end{equation}
Large, certain weights = important features. Small or uncertain weights = weak evidence. Features wit highest $|t|$ values are most significant.

\begin{itemize}
	\item Effect Plot: $effect_j^i = \beta_j \cdot x_j^i$ Contribution of feature j to prediction for instance i.
	\item Different encodings for variables.
\end{itemize}

\subsection{Shrinkage Models}
Using models that keep only the relevant features helps with interpretability.

\subsubsection{Lasso - Least Absolute Shrinkage and Selection Operator}
\begin{itemize}
	\item Adds L1 Norm, Performs feature selection, regularization
	\item Enforces sparsity in Linear model
\end{itemize}
\begin{equation}
	\underset{\beta}{min} \frac{1}{n}\sum:{i=1}^n\lr{y^{(i)}-x_i^T \beta }² + \lambda \norm{\beta}_1
	\label{eq:LASSO-L1}
\end{equation}
As $\lambda$ gets increased, only the most important features remain. This increases interpretability.

\paragraph{About Linear Regression}
\begin{itemize}
	\item Perfectly understood
	\item Only captures linear relationships
	\item Ignores feature interactions, must be added manually ($x_1x_2, x²$ etc.)
	\item Lower accuracy compared to more flexible Models.
\end{itemize}


\subsection{Logistic Regression}
\begin{itemize}
	\item Classification problems with two classes.
	\item Models Probability of one class given the features.
	\item Extension of linear regresion for classification.
	\item Uses the output of a linear regression and maps it to a Probability.
\end{itemize}
\paragraph{Logistic / Sigmoid function}
\begin{equation}
	logistic\lr{\eta} = \frac{1}{1+ \exp{-\eta}}
	\label{eq:Log-sigmoid-fun}
\end{equation}

\paragraph{Logistic Model}
\begin{equation}
	P\lr{y^{(i)} = 1}  = \frac{1}{1 + \exp (- (\beta_0 + \beta_1 x_1^{(i)} + \dots + \beta_p x_p^{(i)}))}
	\label{eq:logisticregression}
\end{equation}
\begin{itemize}
	\item Linear predictor is transformed via the logistic function.
	\item Coeffictions do not act linearly on the probability. Interpretation is happens on the log-odds scale.
	\item A one unit increase in $x_j$ multiplies the odds of $y=1$ vs. $y=0$ by a factor of $\exp(\beta_j)$
	\item Predicts probabilities, which alows for better decision making. Pred with 99\% probability is more confident than one with 51\%.
	\item Relationships between features and outcomes remain transparent.
\end{itemize}


\subsection{Generalized Linear Models}
Linear Regression





\subsection{Generalized Additive Models}



\end{document}




